\documentclass[reqno]{amsart}
%\documentclass[letter,12pt]{report}
%\usepackage{amsaddr}

\usepackage[margin=3cm]{geometry} %set the margins

\usepackage{bm}
\usepackage{hyperref} %[hidelinks] hides the fluorescent boxes and makes all colors black
\hypersetup{
    colorlinks = true,
    linkcolor={gray},
    citecolor={blue!50!black},
    urlcolor={blue!50!black}
}
\usepackage[utf8]{inputenc}
\usepackage{cite}
%\usepackage[sort&compress]{natbib}

\usepackage{amsmath,accents}
%\numberwithin{equation}{section} %numbers the equations according to each section.

\usepackage{siunitx}
\usepackage{mathtools}
\usepackage{amssymb}
\usepackage[normalem]{ulem}
\usepackage{enumitem}
%\graphicspath{{figures/}}

\usepackage{tcolorbox} %allows one to put colored boxes around text and equations (cannot be used with xcolor package}

% For tables
\usepackage{booktabs,siunitx,array,threeparttable}
\sisetup{group-minimum-digits=4}

\def\be{\begin{equation}}
\def\ee{\end{equation}}
\def\e#1{\label{#1}\end{equation}}
\def\bea{\begin{eqnarray}}
\def\eea{\end{eqnarray}}
\def\ea#1{\label{#1}\end{eqnarray}}
\def\bse{\begin{subequations}}
\def\ese{\end{subequations}}
\def\nn{\nonumber}

\def\r#1{(\ref{#1})}
\def\bem#1{\begin{mathletters}\label{#1}}
\def\eml{\end{mathletters}}
\def\df{\Delta\varphi}
\def\agt{\stackrel{>}{_\sim}}
\def\unc#1{\delta \hat{#1}}
\def\uncwrt#1#2{\delta_{\mathrm{#2}} \hat{#1}}
\def\uncwrtpost#1#2#3{\delta_{\mathrm{#2}} \hat{#1}_{(#3)}}
\def\est#1{\hat{#1}}

\newcommand{\ket}[1]{\left| #1 \right>}
\newcommand{\bra}[1]{\left< #1 \right|}
\newcommand{\hham}{\mathcal{H}}
\newcommand{\iu}{{i\mkern1mu}}
\def\braket#1#2{{\langle#1|#2\rangle}}
\def\mean#1{{\langle#1\rangle}}
\newcommand{\myvec}[1]{\accentset{\rightharpoonup}{#1}}
\def\4#1{{\boldsymbol{#1}}}
\def\8#1{\mathbb{E}\left[#1\right]}
\def\func#1#2{\mathrm{#1}\left[#2\right]}
\def\funcwrt#1#2#3{\mathrm{#1}_{\mathrm{#3}}\left[#2\right]}

\DeclareMathSymbol{\mlq}{\mathord}{operators}{``}
\DeclareMathSymbol{\mrq}{\mathord}{operators}{`'}

%\renewcommand\thefootnote{\textcolor{gray}{\arabic{footnote}}} %Not needed if hypersetup is used

\usepackage{titlesec} %titlesec package
\titleformat{\subsection}{\centering\normalfont\bf}{\thesubsection.\quad}{0em}{}{}
\renewcommand{\thesubsection}{\arabic{subsection}}
%\renewcommand{\theequation}{\arabic{section}.\arabic{equation}}

%Allows me to make lists with custom enumeration and also to link a label to that enumeration
\makeatletter
\newcommand{\myitem}[1]{%
\item[#1]\protected@edef\@currentlabel{#1}%
}
\makeatother

\begin{document}

\title{Quantum squeezing cannot beat the standard quantum limit}

\maketitle
\begin{centering}
\author{Liam P. McGuinness}\\
\address{Laser Physics Centre, Research School of Physics, Australian National University, Acton,}\\
\address{Australian Capital Territory 2601, Australia}\\
\email{Email: \href{mailto:liam@grtoet.com}{liam@grtoet.com}}\\
\end{centering}

\begin{abstract}
Quantum entanglement between particles is expected to allow one to perform tasks that would otherwise be impossible \cite{Bell1964,Feynman1982,Deutsch1985, Shor1994, Giovannetti2001}. In quantum sensing and metrology, entanglement is often claimed to enable a precision that cannot be attained with the same number of particles and time, forgoing entanglement \cite{Helstrom1969, Holevo1973a, Giovannetti2004, Giovannetti2006, DemkowiczDobrzanski2012, Zwierz2012, Pezze2018}. % The operative word here is \emph{impossible} -- not technically difficult or hard -- and it sets a high bar to pass. 
Two distinct approaches exist: % in attempts to surpass this precision limit: 
creation of entangled states that either \emph{i)} respond quicker to the signal% (e.g. NOON, CAT or GHZ states)
, or \emph{ii)} are associated with lower noise and uncertainty%, i.e. squeezed states
. Here we show that if our definition of success is -- a precision that is impossible to achieve without entanglement -- then the second approach cannot succeed.\\
\end{abstract}

\section*{Introduction}
%We first give a non-rigorous discussion of measurement precision in the context of spin-1/2 particles and outline arguments towards the proof. % a layperson explanation of the proof.\\
Consider a sensing device composed of $N$ spin-$1/2$ particles which is used to measure some unknown signal $\theta$. %Typical scenarios include measuring the amplitude or direction of a magnetic field, frequency estimation or phase estimation. 
Taking just one spin and using the spin direction $\myvec{S}$ of this particle as a meter, then by error propagation any estimate $\est{\theta}$ of the value of $\theta$ has an uncertainty:\footnote{More fully, we can include the possibility of a finite bias $\epsilon$ and average over $p(\theta = \Theta)$, the \emph{a priori} probability distribution of $\theta$, to obtain the expected uncertainty: $\langle \Delta \est{\theta}\rangle \geq \sqrt{\int^{\Theta} p(\theta = \Theta)\left( \left| \frac{\partial \myvec{S}}{\partial \theta} \right|^{-2} \left(\Delta \myvec{S}\right)^2+\epsilon^2\right) d\theta}$.}:
\be
\Delta \est{\theta} \geq \left| \frac{\partial \myvec{S}}{\partial \theta} \right|^{-1} \Delta \myvec{S}, % + \mathrm{Bias},
\label{eq:SNR}
\ee
where $\Delta \myvec{S}$ is the uncertainty in determining the spin direction. The derivative term is often called the measurement signal, and $\Delta \myvec{S}$ %$(\Delta \myvec{S})^{-1}$ 
the measurement noise\footnote{Incorrectly so, it is the measurement uncertainty not noise.}, so that the `signal-to-noise ratio' determines the measurement precision. With $N$ identical and independent spins, in general one cannot do better than
\be
\Delta \est{\theta} \geq \Delta \est{\theta}_1 / \sqrt{N}, \label{eq:SQL}
\ee
where $\Delta \est{\theta}_1$ bounds the uncertainty using a single particle. Note, for the inequality in Eq.\,\eqref{eq:SQL} to hold, $\Delta \est{\theta}_1$ must represent the best possible uncertainty that can be attained with a single particle, otherwise one could beat the limit simply by improving $\Delta \est{\theta}_1$. Furthermore, in general $\Delta \est{\theta}_1$ is a function of time since the spin response is given by unitary evolution and with more time one can increase the spin response or perform more measurements to improve the uncertainty. Under this definition, Eq.\,\eqref{eq:SQL} is called the standard quantum limit (SQL) and it sets an uncertainty bound per unit time that is impossible to surpass with a given number of identical and independent spins.

There are two approaches to overcoming the SQL using entanglement, either make the spin response greater, or reduce the uncertainty in measuring the spin direction \cite{Gross2012, Ludlow2015, Pezze2018}. The first makes use of entangled NOON, CAT or GHZ states so that, in theory: $\frac{\partial \myvec{S}_N(t)}{\partial \theta} = N \frac{\partial \myvec{S}_1(t)}{\partial \theta}$, whilst $\Delta \myvec{S}_N = \Delta \myvec{S}_1$, where the subscript denotes the number of spins in the sensor \cite{Pezze2018, Gross2012, Ludlow2015, Bollinger1996, Dowling1998, Childs2000, Campos2003, Kok2004, Leibfried2004, Nagata2007, Resch2007, Berry2009, Dorner2009, Escher2011, Giovannetti2011, Daryanoosh2018, Shaniv2018, Xie2021}. In effect a single spin with greater magnetic moment is created such that the response of the entangled device is $N$-fold greater than that of a single spin. It is informative to note here, that $\Delta \myvec{S}$ is the result of sampling from an unknown probability distribution and since both a single entangled state and a single spin provide only one sample, they have the same measurement uncertainty. We will not discuss this approach except to note that when the resources required to generate entanglement are fully accounted for, no improvement over the SQL or $\Delta \est{\theta}_1$ has been demonstrated \cite{McGuinness2021}. The second approach uses entangled squeezed states so that\footnote{Early proposals did not require entanglement to overcome the SQL (see Appendix\,\ref{App:History}). This still leads to confusion as to what constitutes the SQL and whether entanglement is required. For example, photon states with non-Poissonian statistics are also called non-classical states and claimed to surpass the SQL (I prefer the descriptor `non-thermal' since a single unentangled photon or any single particle including a single baseball is a non-classical state in this definition). We explicitly define the SQL as whatever one can achieve without entanglement, thus it is a moot point whether entanglement is required to surpass it.}% The definition of squeezing is also not standardized, sometimes NOON states are called squeezed states \cite{Soerensen2001,Meyer2001}.}
: $\frac{\partial \myvec{S}_N(t)}{\partial \theta} = \frac{\partial \myvec{S}_1(t)}{\partial \theta}$, whereas $\Delta \myvec{S}_N = \Delta \myvec{S}_1/N$. I.e. the spin response of the entangled device remains the same as that of a single spin but the measurement noise reduces by a factor of $N$ \cite{Caves1981, Pezze2018, Gross2012, Ludlow2015, Walls1983, Wodkiewicz1985, Wu1986, Slusher1987, Xiao1987, Polzik1992, Wineland1992, Kitagawa1993, Wineland1994, Sanders1995, Kuzmich1998, Soerensen1998, Brif1999, Kuzmich2000, Meyer2001, Esteve2008, Appel2009, Eberle2010, Gross2010, Koschorreck2010, Koschorreck2010a, Leroux2010, Leroux2010a, LouchetChauvet2010, SchleierSmith2010, Wasilewski2010, Luecke2011, Sewell2012, Aasi2013, Taylor2013, Bohnet2014, Muessel2014, Strobel2014, Kruse2016, Polzik2016, Cox2016, Davis2016, Bohnet2016, Hosten2016, Macri2016, Linnemann2016, Tse2019, Braverman2019, Schulte2020, Malia2020, Bao2020, PedrozoPenafiel2020, Casacio2021, Gilmore2021, Greve2022, Malia2022}.

We can already see that there is a conflict between our explanation of the origin of the measurement uncertainty and what is observed with squeezed states. Shouldn't the uncertainty in estimating the direction of a single squeezed state be the same as for a single spin? Where does this reduced uncertainty come from? Either our layperson description is wrong or squeezed states do not perform as described. We now show that, if entanglement provides any benefit over the SQL, then it must come about from increasing the sensor response to the signal and not through reduced noise\footnote{My apologies for changing between `uncertainty' and `noise', they do not mean the same thing. I would like to consistently use the correct terminology - uncertainty, but the papers in this field refer to noise and I am quoting their terminology. This conflation of terms is the crux of the issue, I encourage the reader to take careful note of this point.}.

\section*{Noise independent quantum precision bound}% Quantum precision bound only depends on detector response to signal.}
For rigour we make two adjustments to our language. First, rather than talk about spin direction, we address the underlying mathematical object -- the state vector. To show that the uncertainty bound obtained from an entangled ensemble containing squeezed states is worse than that of an unentangled ensemble, we use a counting argument based on the number of indivisible state vectors (i.e. a state vector that cannot be further factored into non-separable states) in the ensemble. Henceforth, a state vector, $\ket{\psi}$ always refers to this basic unit we are dealing with -- an indivisible state. A single state vector is never used to describe the quantum state of a separable ensemble, instead we keep note of the number of copies of each state. Secondly, for technical reasons and conciseness we avoid quantitatively defining uncertainty, we define instead the (Fisher) information on $\theta$ denoted $\func{I}{\theta}$, provided by a given state or measurement of that state \cite{Braunstein1994}:
\be
\func{I}{\theta} \equiv \int^{X} \mathrm{d}X \frac{1}{\func{Pr}{X|\theta}}\left(\frac{\partial \func{Pr}{X|\theta}}{\partial \theta}\right)^2, \quad \func{I}{\theta} \equiv \sum_{i=1}^R \frac{1}{\func{Pr}{X_i|\theta}}\left(\frac{\partial \func{Pr}{X_i|\theta}}{\partial \theta}\right)^2
\label{eq:QFI}
\ee
where $\func{Pr}{X|\theta}$ is the conditional probability to obtain the measurement result $X$ given $\theta$ and the (LHS) RHS assumes the measurement outcomes are (continuous) discrete with $R$ possibilities. We note that for any measurement, the estimation uncertainty $\Delta \est{\theta}$ is a monotonically decreasing function of $\func{I}{\theta}$. This observation provides the necessary tools to compare uncertainties, to say one is greater than another, and is sufficient for our purposes.\\

\emph{Key requirement:} The central claim in squeezing enhanced metrology is that squeezed states have an improved intrinsic noise/uncertainty compared to a single spin. It is clear that the response to $\theta$ of the squeezed state, quantified by $\left|\frac{d\ket{\psi(\theta,t)}}{d \theta}\right|$, is not greater than that of a single spin. One reason this must be the case is that if the noise can be reduced by a factor of $N$, then any further improvement in the sensor response would violate the Heisenberg limit.\\% Our proof directly addresses this claim, we show an information bound that just depends on the state response to $\theta$, therefore, if the state response is not increased, the uncertainty cannot be reduced beyond the SQL.

Finally, we only establish a bound on the information provided by a single measurement of the ensemble. Further assumptions for the proof are provided in Appendix\,\ref{App:Assumptions}.\\
%\begin{itemize}
%\item A single measurement of each particle in the ensemble.
%\item Where the path parameterized by $\theta$ is smooth. I.e. the probability distribtion of $\theta$ is smooth.
%\item $\left(\Delta \est{\theta}\right)^2$ denotes the expected mean squared error of the estimator, thus $\Delta \est{\theta}$ is the positive square-root of the expected mean squared error (expected root-mean-square error) of the estimator.
%\item We are calculating the information gain, uncertainty reduction. Or something proportional to this.
%\end{itemize}

%Although it appears like a restrictive list, what types of experiments does this cover?
%Ramsey interferometry, Mach-Zehnder interferometry. Most experimental demonstrations - B field estimation, frequency estimation, phase estimation, operating regime of LIGO (assuming short measurement time so that the gravitational wave signal does not change much).

\emph{Counting approach:} For a measurement device comprised of $N$ spins, assuming we can prove
\begin{quotation}
Statement 1: A single squeezed state does not provide fundamentally more information on $\theta$ than a single spin.
\end{quotation}

It then follows that the information bound on $\theta$ for any $N$ spin ensemble containing $M$ squeezed states is lower than an ensemble of $N$ unentangled spins. The reason being that a squeezed ensemble can be separated into a maximum of $M$ independent state vectors where $M < N$ (this follows from the very definition of entanglement). Assuming these states are independent\footnote{This assumption just ensures that any correlations come solely from entanglement.}, the information provided by the squeezed ensemble is $M$ times that of a single squeezed state and is therefore less than the information provided by the unentangled ensemble. In fact, this counting argument shows that increasing squeezing leads to a worse uncertainty bound since there are less states to average over. Meaning that, if the uncertainty provided by an ensemble containing squeezed states ever surpasses the single particle bound and if the degree of squeezing/entanglement is continuously increased, then at some point the uncertainty must get worse. And the converse: if the measurement uncertainty always improves for arbitrary amounts of squeezing, then the uncertainty never surpasses the single particle bound.

Many mathematical statements equivalent to Statement 1 (excluding the counting argument) have been published (see Appendix\,\ref{App:Proofs}). For example, Wootters showed in 1981 that $\func{I}{\theta}$ can be interpreted as distance metric over quantum states \cite{Wootters1981}, meaning that if the path mapped out by unitary evolution is the same for two states then so is the information. This is equivalent to saying that the uncertainty depends only on the sensor response to $\theta$ \cite{Ou1997, Childs2000, Giovannetti2004, Zwierz2012, Pang2014, Yuan2015, Pang2017, Gorecki2020}, therefore states with no enhanced response provide no fundamental sensing advantage. Braunstein, Caves and Milburn provided most of the mathematical content of the proof by showing that for pure states, $\func{I}{\theta}$ is given solely by the gradient of the state vector with respect to $\theta$, and does not depend on any intrinsic uncertainty of this state \cite{Braunstein1996}. Here we detail the arguments in full.

\section*{Proof - Time-independent multiplicative Hamiltonian: $\hham(\theta) = \theta \cdot \hham$}
Denote a spin-squeezed state as $\ket{\psi_{SS}(\theta,t)}$, and the state of a single spin as $\ket{\psi_1(\theta,t)}$. Denote the maximum information on $\theta$ that can be provided by any measurement on $\ket{\psi_{SS}(\theta,t)}$, $\ket{\psi_1(\theta,t)}$ as $\funcwrt{I}{\theta, t}{SS}$ and $\funcwrt{I}{\theta, t}{1}$ respectively, then we have:\\

\textbf{Claim:} If $\left(\frac{\mathrm{d}\bra{\psi_{SS}(\theta,t)}}{\mathrm{d} \theta}\right)\left(\frac{\mathrm{d}\ket{\psi_{SS}(\theta,t)}}{\mathrm{d} \theta}\right) = \left(\frac{\mathrm{d}\bra{\psi_1(\theta,t)}}{\mathrm{d} \theta}\right)\left(\frac{\mathrm{d}\ket{\psi_1(\theta,t)}}{\mathrm{d} \theta}\right)$ then $\funcwrt{I}{\theta}{SS} \leq \funcwrt{I}{\theta}{1}$.\\

\textbf{Physical interpretation:} Squeezed states are claimed to surpass the SQL by reducing the uncertainty associated with some state, and not by increasing the response to the signal. To refute this claim we need to show that if the gradient of the squeezed state with respect to the signal is the same as that of a single spin, then information bound on the squeezed state is less than or equal to that of a single spin. The problem is that the derivative of a state-vector is also a state vector and we can't say one state vector is greater than another, i.e. $\frac{d\ket{\psi_1}}{d\theta} > \frac{d\ket{\psi_2}}{d\theta}$ is not a mathematical statement. To obtain a non-negative real number, we need to take the inner-product. Since state vectors are normalised, this operation returns the magnitude of the derivative.\\

\textbf{Proof:} Braunstein, Caves and Milburn showed that for a pure state \cite{Braunstein1996}:
\be
\func{I}{\theta} = 4\left[\left(\frac{\mathrm{d}\bra{\psi(\theta,t)}}{\mathrm{d} \theta}\right)\left(\frac{\mathrm{d}\ket{\psi(\theta,t)}}{\mathrm{d} \theta}\right) - \left|\bra{\psi(\theta,t)}\left(\frac{\mathrm{d}\ket{\psi(\theta,t)}}{\mathrm{d} \theta}\right)\right|^2\right].
\label{eq:BCM}
\ee
The time-evolution of an initial state $\ket{\psi_0}$ in response to a Hamiltonian $\hham(\theta,t) = \theta \cdot \hham$, is given by: $\ket{\psi_0} \rightarrow \mathcal{U}(\theta, t)\ket{\psi_0} \equiv \ket{\psi(\theta,t)}$, where $\mathcal{U}(\theta,t) = \func{Exp}{-\iu \theta t \hham /\hbar}$. Working in the Schr\"odinger picture where time-evolution is carried by quantum states, and writing $\ket{\psi_0}$ in terms of the $K$ eigenstates of $\hham$, denoted by their eigenvalues $\ket{\psi_{E_k}}$ with complex amplitude $\alpha_k$: $\ket{\psi_0} = \sum_{k = 1}^m \alpha_i \ket{\psi_{E_k}}$, we have:
\[ \ket{\psi(\theta,t)} = \sum_{k = 1}^K \func{Exp}{-\iu \theta t E_k/\hbar}\alpha_k \ket{\psi_{E_k}}, \]
where the derivative of this state with respect to $\theta$ is:
\[ \frac{\mathrm{d}}{\mathrm{d}\theta}\ket{\psi(\theta,t)} = \sum_{k = 1}^K (-\iu t E_k/\hbar) \func{Exp}{-\iu \theta t E_k/\hbar}\alpha_i \ket{\psi_{E_k}}.\]
For a spin-1/2 with two eigenstates and denoting the eigenvalues of $\hham$ as $\pm \frac{\gamma}{2}$, we have:
\[ \left(\frac{\mathrm{d}\bra{\psi_1(\theta,t)}}{\mathrm{d} \theta}\right)\left(\frac{\mathrm{d}\ket{\psi_1(\theta,t)}}{\mathrm{d} \theta}\right) = \left(\frac{t \gamma}{2\hbar}\right)^2 \left[ |\alpha_1|^2 + |\alpha_2|^2 \right], \]
and
\[ \left|\bra{\psi_1(\theta,t)}\left(\frac{\mathrm{d}\ket{\psi_1(\theta,t)}}{\mathrm{d} \theta}\right)\right|^2 = \left(\frac{t \gamma}{2 \hbar}\right)^2 \left[ |\alpha_1|^2 - |\alpha_2|^2 \right]^2.\]
We can see that $\func{I}{\theta}$ is maximised by initially placing the spin in an equal superposition of eigenstates, so that $|\alpha_1|^2 - |\alpha_2|^2 = 0$. Then: 
%\be
\[\funcwrt{I}{\theta,t}{1} = \left(\frac{t \gamma}{\hbar}\right)^2.\]
%\label{eq:single}
%\ee
We have shown that: $\funcwrt{I}{\theta,t}{1} = 4\left(\frac{\mathrm{d}\bra{\psi_1(\theta,t)}}{\mathrm{d} \theta}\right)\left(\frac{\mathrm{d}\ket{\psi_1(\theta,t)}}{\mathrm{d} \theta}\right)$ and reproduced the well-known Heisenberg limit for a single spin \cite{Bollinger1996, Giovannetti2004, Pang2014}. Inserting into Eq.\,\eqref{eq:BCM}, the equality stated in the claim, the maximum information of the squeezed state is:
\[%\be
\funcwrt{I}{\theta,t}{SS} = \funcwrt{I}{\theta,t}{1} - 4\left|\bra{\psi_{SS}(\theta,t)}\left(\frac{\mathrm{d}\ket{\psi_{SS}(\theta,t)}}{\mathrm{d} \theta}\right)\right|^2.
\]%\label{eq:squeeze} \ee
As the second term is a non-negative number, $\funcwrt{I}{\theta,t}{SS} \leq \funcwrt{I}{\theta,t}{1}$ and the proof is complete.

Whilst the above only considers quantum states, we can explicitly find the measurement associated with $\funcwrt{I}{\theta,t}{1}$. The measurement basis is a projection orthogonal to the eigenstates\footnote{Here orthogonal means at an angle of $90^{\circ}$, not anti-parallel.} \cite{Braunstein1996, Childs2000, Giovannetti2004, Pang2017}, so that if $\hham = \gamma \sigma_z/2$, where $\sigma_z$ is the Pauli-$z$ matrix, a projective measurement in the $x-y$ plane maximises $\func{I}{\theta}$. For a 2-dimensional space with discrete measurement outcomes, we can denote the measurement results as `1' and `0', thus allowing Eq.\,\eqref{eq:QFI} to be expressed as\footnote{We have used $\func{Pr}{0|\theta} = 1 - \func{Pr}{1|\theta}$, for a Bernoulli random variable taking only two values.}:
\be
\func{I}{\theta} = \left|\frac{\partial \func{Pr}{1|\theta}}{\partial \theta}\right|^{2} \frac{1}{\func{Pr}{1|\theta}\left(1-\func{Pr}{1|\theta}\right)}.
\label{eq:prob}
\ee
We can identify $\left|\frac{\partial \func{Pr}{1|\theta}}{\partial \theta}\right|$ and $\sqrt{\func{Pr}{1|\theta}\left(1-\func{Pr}{1|\theta}\right)}$ with the measurement signal and noise in Eq.\,\eqref{eq:SNR}, where the latter is called quantum projection noise \cite{Itano1993}. Note that the noise and signal are not independent, and the description of some quantum states as having intrinsically lower uncertainty is completely absent in this analysis.

In Appendix\,\ref{App:Mixed},\,\ref{App:General} we generalise the proof to the following situations. We show that if we consider:
\begin{itemize}
\item probability mixtures of squeezed states, then the information from these mixed states cannot exceed a single spin.
\item the expected mean information, averaged over the probability distribution $p(\theta = \Theta)$, then squeezed states cannot outperform a single spin.
\item a modified claim concerning the projection of the gradient $\frac{\mathrm{d} \ket{\psi}}{\mathrm{d}\theta}$ orthogonal to $\ket{\psi}$, then a squeezed state cannot cannot outperform a single spin for estimation of signals in arbitrary time-dependent Hamiltonians $\hham(\theta, t)$.
\end{itemize}

\section*{Discussion}
We have proved a powerful and surprising theorem that allows us to immediately exclude any of the methods proposed in \cite{Caves1981, Walls1983, Wodkiewicz1985, Slusher1985, Wu1986, Slusher1987, Xiao1987, Polzik1992, Wineland1992, Kitagawa1993, Wineland1994, Sanders1995, Kuzmich1998, Soerensen1998, Brif1999, Kuzmich2000, Meyer2001, Orzel2001, Esteve2008, Appel2009, Eberle2010, Gross2010, Koschorreck2010, Koschorreck2010a, Leroux2010, Leroux2010a, LouchetChauvet2010, Riedel2010, SchleierSmith2010, Wasilewski2010, Luecke2011, Sewell2012, Aasi2013, Taylor2013, Bohnet2014, Muessel2014, Strobel2014, Kruse2016, Polzik2016, Cox2016, Davis2016, Bohnet2016, Hosten2016, Hosten2016a, Macri2016, Linnemann2016, Tse2019, Braverman2019, Schulte2020, Malia2020, Bao2020, PedrozoPenafiel2020, Casacio2021, Gilmore2021, Colombo2022, Greve2022, Malia2022, Braginsky1974, Braginsky1977, Unruh1979, Braginsky1980, Braginsky1996} from achieving a measurement precision beyond the SQL. Of these works, the following are experimental papers \cite{Xiao1987, Slusher1985, Slusher1987, Polzik1992, Soerensen1998, Kuzmich2000, Meyer2001, Orzel2001, Esteve2008, Appel2009, Eberle2010, Gross2010, Koschorreck2010, Koschorreck2010a, Leroux2010, Leroux2010a, LouchetChauvet2010, Riedel2010, SchleierSmith2010, Wasilewski2010, Luecke2011, Sewell2012, Aasi2013, Taylor2013, Bohnet2014, Muessel2014, Strobel2014, Bohnet2016, Cox2016, Hosten2016, Hosten2016a, Kruse2016, Linnemann2016, Braverman2019, Tse2019, Bao2020, Malia2020, PedrozoPenafiel2020, Casacio2021, Gilmore2021, Colombo2022, Greve2022, Malia2022, Nichol2022}, the majority of which claim to demonstrate a measurement precision that cannot be obtained without entanglement. %\footnote{We need a method to physically characterise or measure the state dependence on $\theta$ since we proved that if the gradient of a spin squeezed state is the same as that of a single state: $\left|\frac{d\ket{\psi(\theta,t)}_{SS}}{d \theta}\right| = \left|\frac{d\ket{\psi(\theta,t)}_1}{d \theta}\right|$, then both states contain the same amount of information on $\theta$. However the gradient of a state-vector is also a vector, not a non-negative real number, so to get a real number we need to take the inner product.  One might argue, isn't it tautological that if the precision is better, then by definition the gradient was increased! Perhaps people mistakenly attribute this to reduced noise and it is just an explanatory error. The problem is that it is generally obvious and easy to check the state dependence on $\theta$. Just look at the measurement results or the expected state evolution, i.e. the quantum phase accumulate per unit time. It is not the case that the phase change per unit time increases and it is mistakenly reported. Furthermore, if this information is not provided, then \emph{de facto}, sufficient information on the measurement uncertainty has not been presented. As an example take the gravitational wave detector LIGO, where squeezed light is inserted into the interferometer. The imparted phase per photon does not increase when squeezed light is inserted into the interferometer. In Ramsey interferometry with squeezed spins, it is also clear that the phase response is not increased, therefore these techniques cannot beat the SQL. A final point which can be subtle to grasp. I am not suggesting that squeezing did not improve the sensitivity of the LIGO interferometer. The results are quite clear in that regard. I am asserting that squeezing did not improve the sensitivity of the LIGO interferometer beyond a strict impossibility limit for the same number of independent input photons and measurement time.}.
This begs the question:
\begin{quotation}
Why are there so many published experimental papers that contradict this proof?
\end{quotation}

The answer is not that our proof is erroneous, but rather that in general, analyses of experiments either misrepresent the achieved precision with respect to the SQL or involve fundamental mistakes. Below we outline some common examples in the literature.\\% We acknowledge that the benefits of entanglement in quantum metrology are so ingrained in the consciousness that it can be difficult to overturn this conception. To help we make some ridiculous and outlandish comparisons, not just because it is fun, but because it can help to illustrate how absurd the claims are.\\

%If someone presents you with a machine that produces more output energy than input, a simple appeal to conservation of energy can be used to correct them, likewise proof of the claim, but it is somewhat non-constructive.  but sometimes it is more informative to go into detail and show exactly where the error in logic occurred.

\noindent\emph{Common reasons to incorrectly claim surpassing the SQL using quantum squeezing:}%The reason seems to be because most experimental papers use a different definition of the SQL or incorrectly report the measurement precision.

\begin{enumerate}
\item \textbf{Comparison to an imperfect limit:} An improvement in measurement precision for a device operating far away from the SQL is demonstrated \cite{Aasi2013, Taylor2013, Tse2019, Casacio2021}. This is not the same as demonstrating a measurement precision beyond the SQL. % If I promise to improve the performance of your car using \emph{Shacky-mo!}, and provide a demonstration where the speed of a run-down old car is improved slightly, you might wonder whether \emph{Shacky-mo!} will work on your car specifically or improve performance in general. In the context of squeezing, we can ask whether the sensitivity improvement is with respect to a fundamental limit that could not have been achieved otherwise. The answer is no\footnote{Take the work of Casacio et.\,al. \cite{Casacio2021}, where squeezed light lowers the microscope noise without increasing the signal and the authors claim to beat the SQL (they use the term shot-noise limit). In response to Referee 1 (in the peer review file) who notes that the actual performance is underwhelming and far from the state-of-the-art, the authors substantially improve the microscope performance by aligning the microscope. If the performance can be improved through alignment it is clear that comparison is not to a strict impossibility limit. In the context of LIGO \cite{Aasi2013, Tse2019} the analogy of a run-down car is less apt, since this is one of the most sensitive detectors that has ever been developed. However, one can clearly read in LIGO's 2022 \href{https://dcc-llo.ligo.org/public/0182/G2200736/001/UpdateonLVKDetectors.pdf}{update} that significant technical noise sources remain which mean that the interferometer is operating far away from the SQL. A factor of 2 -- 3 sensitivity improvement for a device operating orders of magnitude away from the SQL is not the same as demonstrating an improvement beyond this limit.}.

\item \textbf{Tautological definition of the SQL:} The SQL is defined as the precision of a measurement device \emph{as it is and without modifications}. Modifications that utilise entanglement are allowed, whilst forbidding any other changes to the device \cite{Wu1986, Xiao1987}. If the device was already operating at the SQL, and squeezing implemented whilst keeping same number of particles and time, then the sensitivity cannot improve.% on Of course if we forbid anyone changing anything in the measurement device then the precision cannot be improved. The tricky part is to allow modifications that utilise entanglement but forbid any other changes to the device. One example of this definition is the work of \cite{Wu1986, Xiao1987}. 

\item \textbf{Analysis of measurement noise not uncertainty:} Conflation of noise and uncertainty is one of the most common analytical mistakes in spin squeezing, see e.\,g. \cite{Orzel2001, Esteve2008, Appel2009, Gross2010, Koschorreck2010, Koschorreck2010a, Leroux2010, LouchetChauvet2010, Riedel2010, SchleierSmith2010, Sewell2012, Cox2016, Hosten2016, Kruse2016, Colombo2022, Greve2022, Malia2022}. $(\Delta \myvec{S})^2$ is (incorrectly) defined as the measurement variance, such that a variance less than $|S|/4$, where $|S|$ denotes the spin vector length, is evidence of an enhanced precision. To see the error in this logic, allow a projective measurement in the $x$-basis on a single spin pointing in an unknown direction in the $x-y$ plane. The uncertainty in estimating the $x$-component of the spin direction is actually independent of the measurement variance. If ten experiments produce the following data-set: $\{0,0,0,0,0,0,0,0,0,0\}$, the SQL has not been surpassed, even though the variance is less than 1/4. Ultimately we do not want measurements with less \emph{noise}, but measurements that provide lower \emph{uncertainty} on the spin direction. Often prior information from state preparation (and the value of $\theta$) is used to ensure that the measurement statistics have low variance without reducing the uncertainty in estimating $\theta$.

%It is easy to perform measurements with low noise -- throw away the data. Given $N$ identical copies of a state, the noise can be reduced by a factor of $\sqrt{N}$ by measuring only one state. According to squeezing logic, this strategy surpasses the SQL. Thus we can see that some processes which increase the measurement noise actually reduce the measurement uncertainty.

\item \textbf{Replace the signal gradient with the contrast:} The measurement signal (c.f. the derivative in Eqs.\,\eqref{eq:SNR},\,\eqref{eq:prob}), is replaced with a constant term that does not depend on $\theta$ or the measurement basis. When the measurement precision is defined, look for the words: `contrast', `visiblity', `coherence' or `Wineland squeezing parameter' which replace the gradient with the spin vector length \cite{Meyer2001, Esteve2008, Gross2010, Leroux2010, Leroux2010a, Riedel2010, Hosten2016, PedrozoPenafiel2020, Malia2020, Malia2022}. In these works a measurement basis is chosen to reduce the measurement noise without accounting for the commensurate reduction in signal.

%\item \textbf{Incorrect definition of the SQL:} The SQL is commonly defined a time-independent measurement precision or defined with a square-root time dependence, e.g. \cite{Leroux2010a}.
\end{enumerate}

%\textit{There are valid objections to this proof. Only deals with fundamental quantum measurement noise, for other sources of course the noise should be reduced. But entanglement doesn't help with this. Perhaps there is a reason why we can't physically reach this bound for a single spin, and a much stricter limit holds for measurements of a single spin. but we can for an entangled spin. Note for a single spin the optimal measurement is also relatively easy to implement, it is just a projective measurement in the basis spanned by the superposition states. of the spin into one of the eigenstates. Perhaps, decoherence, ...}

\section*{Conclusion}
Ultimately, noise is a statistical sampling phenomenon and reduction of noise through quantum correlations leads to less samples and worse uncertainty. By relating metrological performance to the number of separable states we have proved that squeezed ensembles cannot outperform unentangled ensembles in sensing. The proof was inspired by the proposition of a general quantum mechanical uncertainty limit per wavefunction \cite{McGuinness2021a}, where wavefunctions are considered the fundamental information unit in computers and sensors, not particles. My position is that it this counting argument holds for all entangling procedures, not just squeezing, and applies to general tasks in quantum information like computation. Finally, it should be apparent that in the field of quantum physics, the peer-review process has failed (see also \cite{McGuinness2022, McGuinness2023, McGuinness2023a}) and this has led to a loss of scientific progress.

\bibliographystyle{naturemag}
%\bibliographystyle{unsrt}
\bibliography{}

\section*{Appendix}
\subsection{A confusing history of early ideas to surpass the SQL}\label{App:History}
The idea of using special projective measurements (not states), with reduced noise originated in discussions on improving the sensitivity of gravitational wave detectors \cite{Braginsky1974, Braginsky1977, Unruh1979, Braginsky1996}. These low noise measurements, called quantum non-demolition (QND) measurements, are performed by choosing a measurement operator so that the meter is in an eigenstate of the observable and the variance of the measurement observable goes to zero \cite{Unruh1978, Braginsky1980}. However no entanglement is required to observe such a variance since separable eigenstates also yield zero variance results. One unfortunate relic of this early analysis which persists today is that the SQL is often defined as the minimum obtainable variance on one observable when the variance is equal to that of another conjugate observable. Often this statement is incorrectly reduced to saying that the uncertainty in measuring spin-direction is equivalent to the spin noise, giving the SQL: $\Delta \myvec{S} = \sqrt{|S|}/2$. Thus when the spin noise is $< \sqrt{|S|}/2$, the SQL for measuring spin direction has been surpassed. Initially the idea was to take advantage of better measurements, not entangled states, however around the time of Caves' proposal for squeezed vacuum fluctuations in the context of photon interferometry \cite{Caves1981}, emphasis on states of the meter with reduced uncertainty started. These methods were then extended to spin states in Ramsey interferometers, called spin-squeezed states \cite{Kitagawa1991, Wineland1992, Kitagawa1993}. 

We can summarize some reasons why QND measurements do not surpass the SQL. First, the requirement of perfect prior knowledge about the state of the meter in order to perform a QND measurement is self-defeating\footnote{If we are being pedantic, one can learn about \emph{which} eigenstate the meter is in, and perform discrimination.}. Also there is an inherent lack of recognition of the difference between `uncertainty' and `noise'. It is easy to perform measurements with low noise -- simply throw away data. Given $N$ identical copies of a state, the noise can be reduced by a factor of $\sqrt{N}$ by measuring only one state. In general, QND schemes neglected to consider that the measurement signal and measurement noise/uncertainty are not independent, measurements with reduced noise -- where the meter is near an eigenstate of the measurement operator -- also have less signal. In fact, Wootters showed in 1981 \cite{Wootters1981}, that even in a $2^N$ dimensional Hilbert space the optimal measurement yields a probability function of the form: $\func{Pr}{1|\theta} = \func{cos^2}{\frac{m}{2}(\theta - \theta_0)}$, where $m$ is an integer and $\theta_0$ is a constant; thus when the probability to measure an observable goes to one or zero, the gradient disappears (Unruh used related arguments in response initial proposals \cite{Unruh1978}). Moreso, Itano et. al. noted that since projective measurements are not perfect, it is preferable to choose a measurement that maximises projection noise \cite{Itano1993}. The message here is that simply reducing measurement noise does not guarantee better precision, and in general measurements which maximise noise are better. Finally, Giovannetti, Lloyd and Maccone proved, for a restricted set of Hamiltonians, that entangled states and not measurements are required to surpass the SQL \cite{Giovannetti2006}.

Another relic of early suggestions to use QND measurements to surpass the SQL is their current use in the preparation of ``squeezed" states \cite{Kuzmich2000, Appel2009,   Koschorreck2010, LouchetChauvet2010, SchleierSmith2010, Sewell2012, Bohnet2014, Cox2016, Hosten2016, Hosten2016a, Braverman2019, Malia2020, Greve2022, Malia2022}. The idea is to place an ensemble of spins initialised in $\ket{0}^{\otimes N}$ in an optical cavity, where a $\pi/2$-pulse rotates $\myvec{S}$ into the $x - y$ plane. A QND measurement is then performed on the ensemble to project a sub-ensemble of the spins to either $\ket{0}$ or $\ket{1}$. This measurement is generally a dispersive population measurement of the spins, where an optical probe pulse is inserted into the cavity and the phase (frequency) shift of the output light is recorded. The shift is proportional to the number of atoms projected into $\ket{0},\ket{1}$. Apparently, the ensemble contains an entangled state with less noise and therefore improved measurement precision. If you are bemused as to how to write down the entangled state, then join the club. In these papers you will find a pictorial representation of the quantum state in place of a mathematical description, where a Gaussian blur on a (Bloch) sphere is replaced with an elongated blur on a (Bloch) sphere \cite{Gross2012, Pezze2018}. Note, any state on the Bloch sphere is isomorphic to a single spin, and is generally represented by an arrow in the spin direction, not a blur. 

After the QND measurement, I would describe the ensemble as being in a probability mixture of $\ket{0}, \ket{1}$ and $\ket{x}$. To refute this description at least one piece of information is required: a definition of the phase relationship between entangled basis states. For example, one can write an entangled Bell state as: $(\ket{00} + \ket{11})/\sqrt{2}$, and although it is generally overlooked, a well-defined phase between the states $\ket{00}$ and $\ket{11}$ is required to generate and make use of this entangled state (in this example the state can be written $(\ket{00} + e^{\iu \phi }\ket{11})/\sqrt{2}$ with $\phi = 0$). Otherwise we have a mixed state and the best we can say is that for any measurement we will record either a `0' or a `1', each with 50\% probability. None of the papers that use QND preparation of a ``squeezed" state define or make use of any phase relationship between the entangled basis states. It is true that knowledge of the phase of the initial $\pi/2$-pulse and the intensity the probe light is required when measuring these states, but these parameters do not determine $\phi$.\\

\subsection{Proof assumptions}\label{App:Assumptions}
For completeness, we list some assumptions of the proof below. In particular we assume:
\begin{itemize}
\item The postulates of quantum mechanics.
\item The entire dependence on $\theta$ is encoded in $\hham(\theta,t)$, the starting state $\ket{\psi_0}$ and the measurement operator are independent of $\theta$.
\item Apart from the initial state vectors that evolve to $\ket{\psi_{SS}(\theta,t)}, \ket{\psi_1(\theta,t)}$, everything else in the external universe is the same when comparing unentangled and entangled ensembles. The prior information on $\theta$ and the ability to apply control fields is the same. In fact, the squeezing community claim that even unitary evolution is the same for these two states, therefore the only difference of consequence is the starting state.
\item Each copy of the same state is independent and measurements on copies of the same state are identical. This assumption just ensures that any enhancement over the SQL comes solely from entanglement, not from other correlations, additional information or different control on copies.\\
\end{itemize}

\subsection{Alternative proofs of a noise independent uncertainty limit}\label{App:Proofs}
From my perspective, the number of published works available highlights how big the blind-spot in the physics community is. Indeed the mathematics is well-known, it seems like the main difficulty is entertaining the idea that entanglement could lead to a worse measurement uncertainty, or having the conviction to follow the mathematics despite being told otherwise. The following gives some sketches for alternative approaches to obtaining the same proof as in the main text. They are not rigorous and more work is required to close technical details. 
\begin{enumerate}
\item We could make the observation that most experimental implementations of squeezing such as quadrature detection in optical interferometers or spin readout in Ramsey interferometers result in a binary outcome measurement. Thus we can use the expression given in Eq.\,\eqref{eq:prob}:
\[\func{I}{\theta} = \left|\frac{\partial \func{Pr}{1|\theta}}{\partial \theta}\right|^{2}\frac{1}{\func{Pr}{1|\theta}\left(1-\func{Pr}{1|\theta}\right)}.\]
It is left as an exercise to show that any modifications to $\func{Pr}{1|\theta}$ that increase $\func{I}{\theta}$ necessarily increase $\left|\frac{\partial \func{Pr}{1|\theta}}{\partial \theta}\right|$. More involved is to show that this is the case even for probability distributions taking $2^N$ discrete values, thus we do not need to restrict ourselves to a two-dimensional probability space.

\item Instead of operating in the Schr\"odinger picture we could analyse how unitary evolution acts as a mapping of quantum states in the Heisenberg picture. Then, for $\hham(\theta, t) = \theta \hham$, we have : $\frac{\mathrm{d} \mathcal{U}(\theta,t)\ket{\psi_0}}{\mathrm{d}\theta} = -\iu t \hham \, \mathcal{U}(\theta, t)/\hbar \ket{\psi_0}$. As $\mathcal{U}$ and $\hham$ commute, expressing $\ket{\psi_0}$ in terms of the eigenstates of $\hham$, we have:
%\begin*{align}
\[
\begin{split}
& \left(\frac{\mathrm{d}\bra{\psi(\theta,t)}}{\mathrm{d} \theta}\right)\left(\frac{\mathrm{d}\ket{\psi(\theta,t)}}{\mathrm{d} \theta}\right) = \bra{\psi_0}\left(-\iu t \hham \, \mathcal{U}(\theta,t)/\hbar\right)^{\dagger}\left(-\iu t \hham \, \mathcal{U}(\theta,t)/\hbar\right)\ket{\psi_0} \\
& = (t/\hbar)^2 \sum_{k = 1}^K \bra{\psi_{E_k}} \left(\func{Exp}{\iu \theta t E_k/\hbar}E_k \alpha_i^{\ast}\right)\left(\alpha_i E_k\func{Exp}{-\iu \theta t E_k/\hbar}\right) \ket{\psi_{E_k}}\\
& = (t/\hbar)^2 \sum_{k = 1}^K  E_k^2 |\alpha_i|^2.\\
& \left|\bra{\psi(\theta,t)}\left(\frac{\mathrm{d}\ket{\psi(\theta,t)}}{\mathrm{d} \theta}\right)\right|^2 = \left|-\iu t/\hbar \bra{\psi_0}\mathcal{U}^{\dagger} \hham \,\mathcal{U}\ket{\psi_0}\right|^2 = (t/\hbar)^2\left(\sum_{k = 1}^K  E_k |\alpha_i|^2\right)^2.
\end{split}
\]
%\end*{align}
Thus maximising information on $\theta$ means maximising the variance of $\hham$\footnote{In general we want to maximise the variance of the generator of translation with respect to $\theta$.}: $(\Delta \hham)^2 \equiv \langle \hham^2 \rangle - \langle \hham \rangle^2$ with respect to $\ket{\psi_0}$. For unitary evolution this also means maximising the variance of $\hham$ on the output state. We have:
\[ \func{I}{\theta} = 4\left(\Delta \hham\right)^2.\]
Due to the correspondence between different pictures we can see that increasing $\left(\Delta \hham\right)^2$ in this context only comes about through increasing the state response to $\theta$.

\item We can follow Wootters \cite{Wootters1981} and show that $\func{I}{\theta}$ is a metric on state-vectors, characterised by the distance between the start and final state. As this distance is parameterized by the state evolution it follows that if state evolution is the same for two states, then the information is also the same. Equivalently, the information provided by a path of states is given by the gradient vector. We also need to use the fact that state vectors are normalised, so they have the same length.

\item In Refs. \cite{Wootters1981, Braunstein1994, Childs2000, Giovannetti2006, Jones2010, Pang2017} amongst others, an uncertainty bound depending only on the maximal difference in eigenvalues of $\frac{\partial \hham(\theta,t)}{\partial \theta}$ is provided. Minimising the uncertainty requires placing the meter in an equal superposition of states with extremal eigenvalues, $\mu_{\mathrm{max}}(t), \; \mu_{\mathrm{min}}(t)$:
\[
\ket{\psi(t)} = \frac{1}{\sqrt{2}}\left(\ket{\mu_{\mathrm{max}}(t)} + e^{\iu \Phi}\ket{\mu_{\mathrm{min}}(t)}\right),
\]
where $\Phi$ is an arbitrary phase. The optimal state is isomorphic to a single spin, since it is just a 2-dimensional space and the measurement uncertainty per unit time is limited by the maximum response speed of the meter to changes in $\theta$.

\item We could observe that the uncertainty on $\theta$ in Eq.\,\eqref{eq:SNR} is neatly divided into two time-zones (this is in contrast to Eq.\,\eqref{eq:prob}). The first contribution is during sensor interaction with the signal and the second occurs after interaction and is just readout of the sensor. Now the benefit from squeezing closely resembles the original proposals to surpass the SQL (see Appendix\,\ref{App:History}). But Giovannetti, Lloyd and Maccone proved that entanglement cannot provide any benefit at the measurement stage for a restricted class of Hamiltonians \cite{Giovannetti2006}. For a rigorous proof we also need to show that the two terms in Eq.\,\eqref{eq:SNR} are independent. The temporal separation achieves this to the extent required, since after interaction, the first term is fixed and we do not yet need to make any decision on the measurement.

\item We could apply results from quantum state tomography, where the uncertainty in estimating the spin direction of an unknown state just depends on the number of copies of the state \cite{ODonnell2016, Aaronson2018}.

\item We could observe the hint provided by looking at state discrimination, i.e. when we know that $\theta$ can take only one of two discrete values. Childs, Renes and Preskill \cite{Childs2000} showed that driving the state to one of two orthogonal states minimises the discrimination uncertainty and the best one can do is to drive as quickly as possible. There is no benefit to reducing readout noise here since it already goes to zero.\\
\end{enumerate}

\subsection{Proof for average expected information and mixed states}\label{App:Mixed}
The key to these extended proofs is that we are dealing with probability distributions over non-negative real numbers, so we don't have to worry about any complex number magic. In comparing summations of non-negative real numbers where each term in one sum is less than each corresponding term in another sum, by convexity we have an inequality.\\

\noindent \textbf{Mixed states:} To prove the claim when the squeezed state is a mixed state $\ket{\rho_{0,SS}}$, we can write the initial state as probability distribution over $L$ pure states, with probabilities $p_l$: $\ket{\rho_{0,SS}} = \sum_{l=1}^L p_l \ket{\psi_{0,SS}}_l$. As we have assumed that the initial state does not depend on $\theta$, we can use the property that the Fisher information is convex over mixed states when $p_l$ is independent of $\theta$ \cite{Liu2020}.\\

%Then as the information from any individual pure state does not exceed a single spin, the average information from the mixed state cannot either. This is equivalent to  Although, here we do not need to assume independence of probabilities because the optimal state independent of $p_l$.

%For mixed states, when the probabilities are parameter independent, i.e. when $\func{Pr}{\theta} = \mathrm{Pr}$, then the quantum Fisher information is convex:
%\[ \func{I}{\func{Pr}{\cdot}\rho_1(\theta)+(1-\func{Pr}{\cdot})\rho_2(\theta),H} \leq \func{Pr}{\cdot}\func{I}{\rho_1(\theta),H}+(1-\func{Pr}{\cdot})\func{I}{\rho_2(\theta),H}\]


\noindent\textbf{Average expected information:} Eq.\,\eqref{eq:QFI} just provides the information around a given value of $\theta$. The expected information, averaged over the prior probability distribution of $\theta$ is:
\be \langle \func{I}{\theta}\rangle = \int^{\Theta} \mathrm{d}\Theta p(\theta=\Theta) \func{I}{\Theta}, \quad \langle \func{I}{\theta}\rangle = \sum_{j=1}^J p(\theta = \Theta_j) \func{I}{\Theta_j}
\ee
where $p(\theta = \Theta)$ is the prior probability that $\theta$ takes on some value $\Theta$. The LHS expression is the expected information when the prior probability is continuous and the RHS when $\theta$ can take $J$ discrete values.

Given that the information on any value of $\theta$ is not greater than a single spin, we can show that this holds for the expected information averaged over the probability distribution of $\theta$. First we note that the optimal state is independent of the value of $\theta$ and therefore also independent of the prior probability of $\theta$. It is not however the case that the optimal measurement is independent of prior probability distribution of $\theta$. To additionally prove that no measurement can produce more information, we use the result of Giovannetti, Lloyd and Maccone who showed that entanglement does not provide any advantage at the measurement stage \cite{Giovannetti2006}. Thus we only need to consider the amount of information encoded in the state (before measurement).\\

\subsection{Proof - Arbitrary time-dependent Hamiltonian: $\hham(\theta, t)$}\label{App:General}
We motivate the proof by noting that if we could solve the unitary evolution of a single spin and show that it remains in the optimal state for all evolution times, then this would imply $\left|\bra{\psi_1(\theta,t)}\left(\frac{\mathrm{d}\ket{\psi_1(\theta,t)}}{\mathrm{d} \theta}\right)\right|^2 = 0$. However, without a general method of solving time-evolution, even when restricting ourselves to a single spin, we cannot make that statement. An equivalent approach is to use the information bound proved by Pang and Jordan for general time-dependent Hamiltonians \cite{Pang2017}:
\be
\func{I}{\theta,t} \leq \left| \int^t_{0} \mathrm{d}t' \left(\mu_{\mathrm{max}}(t') - \mu_{\mathrm{min}}(t')\right)/\hbar \,\right|^{2},
\label{eq:Pang}
\ee
where $\mu_{\mathrm{max}}(t), \; \mu_{\mathrm{min}}(t)$ are the maximum, minimum eigenvalues of $\frac{\partial \hham(\theta,t)}{\partial \theta}$. Eq.\,\eqref{eq:Pang} expresses a simple concept, the information per unit time is limited by the maximum response speed of the meter to changes in $\theta$. This quantity is characterised by the eigenvalues of $\frac{\partial \hham(\theta,t)}{\partial \theta}$ (c.f. how the eigenvalues of $\hham$ determine the evolution speed of any quantum state). Minimising the uncertainty requires placing the meter in an equal superposition of states with extremal eigenvalues, $\mu_{\mathrm{max}}(t), \; \mu_{\mathrm{min}}(t)$:
\[
\ket{\psi^{\mathrm{Opt}}(t)} = \frac{1}{\sqrt{2}}\left(\ket{\mu_{\mathrm{max}}(t)} + e^{\iu \Phi}\ket{\mu_{\mathrm{min}}(t)}\right),
\]
where $\Phi$ is an arbitrary phase \cite{Braunstein1996, Giovannetti2006, Childs2000}.

If we could solve the unitary evolution under $\hham(\theta, t)$ and show that $\ket{\psi_1(\theta,t)} = \ket{\psi^{\mathrm{Opt}}(t)}$, (and find a projective measurement that saturates the bound) then we can saturate Eq.\,\eqref{eq:Pang} and show that $\funcwrt{I}{\theta,t}{1} = 4\left(\frac{\mathrm{d}\bra{\psi_1(\theta,t)}}{\mathrm{d} \theta}\right)\left(\frac{\mathrm{d}\ket{\psi_1(\theta,t)}}{\mathrm{d} \theta}\right)$. Unfortunately in general this is not possible, for example when both the eigenvalues and eigenstates of $\hham$ depend on $\theta$. Therefore we need to modify our claim to address not the gradient of $\ket{\psi}$, but rather the projection of $\frac{\mathrm{d}\ket{\psi}}{\mathrm{d} \theta}$ orthogonal to $\ket{\psi}$: $\left(\frac{\mathrm{d}\ket{\psi}}{\mathrm{d} \theta}\right)_{\perp} \equiv \frac{\mathrm{d}\ket{\psi}}{\mathrm{d}\theta}-\ket{\psi}\bra{\psi}\left(\frac{\mathrm{d}\ket{\psi}}{\mathrm{d}\theta}\right)$.\\

\textbf{Claim:} If $\left(\frac{\mathrm{d}\bra{\psi_{SS}(\theta,t)}}{\mathrm{d} \theta}\right)_{\perp}\left(\frac{\mathrm{d}\ket{\psi_{SS}(\theta,t)}}{\mathrm{d} \theta}\right)_{\perp} = \left(\frac{\mathrm{d}\bra{\psi_1(\theta,t)}}{\mathrm{d} \theta}\right)_{\perp}\left(\frac{\mathrm{d}\ket{\psi_1(\theta,t)}}{\mathrm{d} \theta}\right)_{\perp}$ then $\funcwrt{I}{\theta}{SS} = \funcwrt{I}{\theta}{1}$.\\

\textbf{Proof:} As noted by Braunstein, Caves and Milburn \cite{Braunstein1996}
\[ \left(\frac{\mathrm{d}\bra{\psi(\theta,t)}}{\mathrm{d} \theta}\right)_{\perp}\left(\frac{\mathrm{d}\ket{\psi(\theta,t)}}{\mathrm{d} \theta}\right)_{\perp} = \left(\frac{\mathrm{d}\bra{\psi(\theta,t)}}{\mathrm{d} \theta}\right)\left(\frac{\mathrm{d}\ket{\psi(\theta,t)}}{\mathrm{d} \theta}\right) - \left|\bra{\psi(\theta,t)}\left(\frac{\mathrm{d}\ket{\psi(\theta,t)}}{\mathrm{d} \theta}\right)\right|^2 \]
and Eq.\,\eqref{eq:BCM} proves the claim. This statement is equivalent to looking at the data-set $\{X_i|\theta\}$ resulting from projective measurements of $\ket{\psi_{SS}(\theta,t)}$. If this data-set shows no greater dependence on $\theta$ than a single spin, then it cannot contain fundamentally more information on $\theta$ than projective measurements of a single spin.\\

%\subsection*{Philosophical question}
%Who should claim the proof? Clearly it was shown by Braunstein, Caves and Milburn. This is more a discussion on what constitutes a proof. In mathematics, given some axioms, it is generally fixed what is provable or not provable in that axiomatic system. So one can say that the person writing the axioms has proved everything it is just a crank the handle type operation. But in general, some problems are described to require some special insight in order to solve them. The person providing that insight is usually credited with the proof, and sometimes the subsequent people who provide limited insight but do the arduous crank-the-handle work in extending to other situations or rigorously checking and defining the proof are also credited. But how do we quantify insight? It is ill-defined and a grey area. One person, could claim that it is obvious to them and they proved it, but did not provide the work, in other cases this is not so.

%For a mathematical journal, field it would be obvious that BCM did the mathematical proof. However in the past 30 years since this work, it is also obvious that they (and subsequent people) did not recognize the implications of this proof. Grigori Perelman 

\subsection{Comment on vacuum noise in an interferometer}
It is not lost on me that this proof seems to contradict fundamental descriptions of interferometer performance. This section tries to reduce some of the cognitive dissonance involved in accepting the proof.\\

In his seminal paper on quantum mechanical noise in an interferometer, Caves noted a fundamental noise source that ``can be attributed to vacuum (zero-point) fluctuations in the electromagnetic field, which enter the interferometer from the unused input port" \cite{Caves1981}. It is commonly understood that this noise cannot be removed, except through inserting squeezed vacuum into the unused input port. There are in fact alternative descriptions of the noise which are equally valid. Caves noted an equivalent view that attributes radiation-pressure fluctuations to the ``random scattering of the input photons at the beam-splitter". Under this perspective, when $N$ photons of perfectly squeezed light are incident on the beam-splitter, they all distribute to the same arm of the interferometer and exit at the same port. Importantly the phase shift experienced by these photons is the same as for a single photon. In effect the squeezed photons collectively act as a single photon, so that the photon statistics and information on the phase shift is identical to what would be observed by passing a single photon through the interferometer.

But we can go one step further and implement a technique to block the unused input port of the interferometer and prevent vacuum fluctuations from entering. To do this, we need to temporally separate photons in one arm of the interferometer from photons in the other arm. To illustrate the idea, we use photons produced from a pulsed laser, and replace the beam-splitter with a flippable (double-sided) mirror, starting flipped down so that incident photons enter the longitudinal arm of the interferometer. The first pulse sends photons into this arm after which the mirror is flipped up, so that photons in the second pulse go into the transverse arm. The mirror remains in this position so that photons returning from the longitudinal arm are sent out the antisymmetric port of the interferometer, after which it is flipped down allowing photons returning from the transverse arm to output at the same port. If we can replace the photodiode normally positioned at the antisymmetric output of the interferometer with a phase sensitive detector, this would allow us to compare the phases of both optical pulses and detect a shift in one arm of the detector.

The fact that photodiodes are sensitive only to the light intensity and not the phase would seem to forbid this mode of operation. However coherent ensembles of spins are phase sensitive detectors. The idea is that the first (longitudinal) optical pulse performs a $\pi/2$-rotation on the ensemble, where the optical phase is stored coherently in the spin direction of the ensemble. This phase can be `readout' with the second (transverse) optical pulse which performs a $\pi/2$-rotation around an axis determined by the relative phase of the two pulses (see \cite{McGuinness2021} for a demonstration with microwave photons). If the energy level of the spins is perfectly resonant to the laser frequency, time-delays between the pulses have no effect on the final state of the spins. In fact we do not even need an interferometer \emph{per se}, to detect gravitational waves. We can remove the mirror and make use of just the longitudinal arm of the interfermeter. Now the phase of returning photons is simply compared to the phase of photons outputted from the laser directly which are either simultaneously or sequentially applied to an ensemble of spins. This mode of operation only allows detection of gravitational waves along one axis and requires that the laser has good phase stability.

This perspective is useful in that it demonstrates that vacuum fluctuations can be removed from an interferometer used to detect gravitational waves. Of course, they are replaced with something else here -- quantum projection noise when the spin state of the ensemble is readout -- thus we do not win overall. The benefit of this second view, to me, is that does not produce the same temptation to reduce the noise in the interferometer. The problem is reduced to estimating the unknown direction of an ensemble of spins. While not as alluring, this perspective better aligns intuition with the correct results.

\end{document}